% ============ 第4章 §4.4 Hilbert 变换的 L^p 有界性 (原书 pp.89-90) ============

\section{\texorpdfstring{Hilbert 变换的 $L^{p}$ 有界性}{Hilbert 变换的 Lp 有界性}}\label{sec:4.4}

这里, 利用 $L^{p}$ 空间的对偶性（定理~\ref{thm:2.19}）
\[
\|f\|_{L^{p}} = \sup\left\{\|f \cdot g\|_{L^{1}} : g \in L^{p'}(\mathbb{R}), \|g\|_{L^{p'}} = 1\right\},
\]
与 $H$ 在 $L^{2}(\mathbb{R})$ 有界, 以及对一切 $\lambda > 0$, $f \in L^{1}(\mathbb{R}) \cap L^{2}(\mathbb{R})$, $\lambda|\{|Hf| > \lambda\}| \lesssim \|f\|_{L^{1}}$, 以及 $f, g \in L^{2}(\mathbb{R})$ 时,
\[
\int_{\mathbb{R}}f(x)Hg(x)\dd x = -\int_{\mathbb{R}}g(x)Hf(x)\dd x,
\]
来考察 Hilbert 变换的 $L^{p}$ 有界性. 虽然是用极限记号定义的, 但实质上 $H$ 是卷积, 所以这个等式成立. 即, 目标是: 对 $f \in L^{p}(\mathbb{R}) \cap L^{2}(\mathbb{R})$ 与 $1 < p < \infty$,
\[
\|Hf\|_{L^{p}} \lesssim \|f\|_{L^{p}}
\]
成立. 这里, $\lesssim$ 中包含的常数不依赖于 $f$, 只依赖于 $p$. 如果能做到这一点, 就能由稠密性将 $H$ 定义为 $L^{p}(\mathbb{R})$ 上的有界线性算子.

考虑 $\chi_{[-1,1]}$ 可知 Hilbert 变换不是 $L^{1}$ 有界的, 但证明对其他的 $p \in (1, \infty)$, Hilbert 变换是 $L^{p}$ 有界的. 重要的是, 从这以后, 除上述对偶公式外, 并不特别需要 $H$ 的结构. 下面的命题称为弱 (1,1) 与 $L^{p}$ 有界性的插值, 即 Marcinkiewicz 插值定理, 它不依赖于维数, 所以作为一般维数的命题来证明.

\begin{theorem}[Marcinkiewicz 插值定理]\label{thm:4.5}
设 $1 < q < p < \infty$. 算子 $S: L^{p}(\mathbb{R}^{n}) \to L^{p}(\mathbb{R}^{n})$ 满足下列条件.

(1) $f \in L^{p}(\mathbb{R}^{n})$ 时, $Sf$ 是非负值函数.

(2)（齐次性）$a \in \mathbb{C}$, $f \in L^{p}(\mathbb{R}^{n})$ 时, 对几乎所有的 $x \in \mathbb{R}^{n}$, $S[af](x) = |a|Sf(x)$ 成立.

(3)（次可加性）$f, g \in L^{p}(\mathbb{R}^{n})$ 时, 对几乎所有的 $x \in \mathbb{R}^{n}$, $S[f+g](x) \leq Sf(x) + Sg(x)$ 成立.

(4)（弱 (1,1) 有界性）对一切 $f \in L^{1}(\mathbb{R}^{n}) \cap L^{p}(\mathbb{R}^{n})$ 与 $\lambda > 0$, $|\{Sf > \lambda\}| \lesssim \frac{1}{\lambda}\|f\|_{L^{1}}$ 成立.

(5)（$L^{p}$ 有界性）对一切 $f \in L^{p}(\mathbb{R}^{n})$, $\|Sf\|_{L^{p}} \lesssim \|f\|_{L^{p}}$ 成立.

此时, 对一切 $f \in L^{p}(\mathbb{R}^{n}) \cap L^{q}(\mathbb{R}^{n})$, $\|Sf\|_{L^{q}} \lesssim \|f\|_{L^{q}}$ 成立. 不过, $\lesssim$ 中包含的常数只依赖于弱 (1,1) 有界性的常数, $L^{p}$ 有界性的常数, $p$, $q$.
\end{theorem}

作为这样的 $S$ 的例子, 除极大算子 $M$ 之外, 还有 Hilbert 变换 $H$ 与绝对值复合而成的 $f \mapsto |Hf|$. 把定理的叙述和证明与定理~\ref{thm:1.3} 比较也是一件乐事.

\begin{proof}
用定理~\ref{thm:3.3} 改写 $L^{q}$ 范数.
\begin{equation}\tag{4.30}
\|Sf\|_{L^{q}}^{q} = q2^{q}\int_{0}^{\infty}\lambda^{q-1}|\{Sf > 2\lambda\}|\dd\lambda.
\label{eq:4.30}
\end{equation}
这里, 对 $\lambda > 0$, 由次可加性, 弱 (1,1) 有界性, $L^{p}$ 有界性以及 Chebyshev 不等式
\[
|\{Sf > \lambda\}| \lesssim \frac{1}{\lambda^{p}}\int_{\mathbb{R}^{n}}Sf(x)^{p}\dd x,
\]
得
\begin{align*}
|\{Sf > 2\lambda\}| &\leq |\{S[f\chi_{\{|f| > \lambda\}}] > \lambda\}| + |\{S[f\chi_{\{|f| \leq \lambda\}}] > \lambda\}| \\
&\lesssim \frac{1}{\lambda}\int_{\{|f| > \lambda\}}|f(x)|\dd x + \frac{1}{\lambda^{p}}\int_{\{|f| \leq \lambda\}}|f(x)|^{p}\dd x \\
&\lesssim \int_{\mathbb{R}^{n}}\left(\frac{\chi_{\{|f| > \lambda\}}(x)|f(x)|}{\lambda} + \frac{\chi_{\{|f| \leq \lambda\}}(x)|f(x)|^{p}}{\lambda^{p}}\right)\dd x.
\end{align*}
因此, 对此不等式乘以 $\lambda^{q-1}$ 再积分,
\begin{align*}
\|Sf\|_{L^{q}} &\lesssim \int_{0}^{\infty}\left\{\int_{\mathbb{R}^{n}}\left(\frac{\chi_{\{|f| > \lambda\}}(x)|f(x)|}{\lambda^{2-q}} + \frac{\chi_{\{|f| \leq \lambda\}}(x)|f(x)|^{p}}{\lambda^{p-q+1}}\right)\dd x\right\}\dd\lambda
\end{align*}
成立. 用 Fubini 定理交换积分顺序后, 关于 $\lambda$ 积分, 得 $\|Sf\|_{L^{q}} \lesssim \|f\|_{L^{q}}$.
\end{proof}

\begin{theorem}[Hilbert 变换的 $L^{p}$ 有界性]\label{thm:4.6}
对一切 $1 < p < \infty$, 存在满足关系式 $c_{p} = c_{p'}$ 的常数族 $\{c_{p}\}_{1 < p < \infty}$, 使得对一切 $f \in L^{2}(\mathbb{R}) \cap L^{p}(\mathbb{R})$, $\|Hf\|_{L^{p}} \leq c_{p}\|f\|_{L^{p}}$ 成立.
\end{theorem}

注意, $p'$ 的共轭指数是 $p$, 且 $p \in (1, 2)$ 与 $p' \in (2, \infty)$ 等价, 所以对一切 $1 < p < \infty$ 关系式 $c_{p} = c_{p'}$ 成立与对一切 $1 < p < 2$ 关系式 $c_{p} = c_{p'}$ 成立是等价的.

\begin{proof}
分别证明 $1 < p < 2$ 的情形与 $2 < p < \infty$ 的情形即可. $1 < p < 2$ 的情形用定理~\ref{thm:4.5} 插值定理~\ref{thm:4.4} 与 $L^{2}$ 有界性. 如此, 对 $1 < p \leq 2$, 对一切 $f \in L^{2}(\mathbb{R}) \cap L^{p}(\mathbb{R})$, $\|Hf\|_{L^{p}} \leq c_{p}\|f\|_{L^{p}}$ 成立. 对 $2 < p < \infty$, 令 $c_{p'} = c_{p}$, 用对偶性证明对 $2 < p < \infty$, 对一切 $f \in L^{2}(\mathbb{R}) \cap L^{p}(\mathbb{R})$, $\|Hf\|_{L^{p}} \leq c_{p}\|f\|_{L^{p}}$ 也成立. 具体地如下计算. 任意固定 $\varepsilon > 0$, 取满足
\begin{equation}\tag{4.31}
\|Hf\|_{L^{p}} \leq (1+\varepsilon)\left|\int_{\mathbb{R}}Hf(x) \cdot \overline{g(x)}\dd x\right|, \quad \|g\|_{L^{p'}} = 1
\label{eq:4.31}
\end{equation}
的 $g \in L^{2}(\mathbb{R}) \cap L^{p'}(\mathbb{R})$. 由于 $f, g \in L^{2}(\mathbb{R})$, 满足 $\int_{\mathbb{R}}Hf(x) \cdot \overline{g(x)}\dd x = -\int_{\mathbb{R}}f(x) \cdot \overline{Hg(x)}\dd x$. 用 Hölder 不等式, 得 $\|Hf\|_{L^{p}} \leq (1+\varepsilon)\|f\|_{L^{p}} \cdot \|Hg\|_{L^{p'}}$. 由 $1 < p' < 2$, 得 $\|Hf\|_{L^{p}} \leq c_{p'}(1+\varepsilon)\|f\|_{L^{p}}$. 由于 $\varepsilon > 0$ 任意, 得 $\|Hf\|_{L^{p}} \leq c_{p'}\|f\|_{L^{p}}$. 因此, 定理~\ref{thm:4.6} 证明完毕.
\end{proof}

\begin{exercise}\label{exer:4.11}
设 $1 \leq p < \infty$, 证明 $L^{p}(\mathbb{R}) \cap L^{2}(\mathbb{R})$ 在 $L^{p}(\mathbb{R})$ 中稠密.
\end{exercise}
